\section{The complementary Bob angles}\label{app:bob}
For Bob testing angles outside the contact cap, two positive source
coefficients control the negative state square directly. This gives a
second fixed certificate on $|v|\geq1/16$, for every Alice angle.

Use the compact testing coordinates
\[
 (a,b)=\frac{(u,1)}{\sqrt{1+u^2}},\qquad
 (c,e)=\frac{(v,1)}{\sqrt{1+v^2}},\qquad
 a^2+b^2=c^2+e^2=1.
\]
They include the projective endpoints by continuity. In the original
state coordinates of \eqref{eq:slack}, put
\begin{gather*}
 E_B=aW-bY,\quad K_B=aX-bZ,\quad J_B=eE_B+cK_B,\quad U_B=eZ-cY,\\
 E_\lambda=eW+cX,\qquad L_A=bX+aZ,\qquad
 M_B=b(cW-eX)+a(cY-eZ).
\end{gather*}
Then the Bell slack has the compact expression
\begin{align}
 S_c={}&\lambda E_\lambda^2+\beta L_A^2+nX^2+mM_B^2+CE_B^2+AY^2
 +(D+d/2)J_B^2\notag\\
 &+(\alpha+d/2)U_B^2+\tfrac d2(L_A^2-Z^2).
 \label{eq:bobparent}
\end{align}
The identity $L_A^2-Z^2=X^2-K_B^2$ leaves one negative square to absorb.
Because $J_B-eE_B=cK_B$, the terms with coefficients $C$ and $D$
give exactly
\begin{align}
 CE_B^2+DJ_B^2-\tfrac d2K_B^2
 &=C(cE_B)^2+(eE_B,J_B)S_B(eE_B,J_B)^T\notag\\
 &\quad+\tau_B(c^2-1/257)K_B^2,\qquad
 S_B=\begin{pmatrix}C-\tau_B&\tau_B\\\tau_B&D-\tau_B\end{pmatrix},
 \quad \tau_B=257d/2.
 \label{eq:boballocation}
\end{align}
The last square has nonnegative weight precisely when $|v|\geq1/16$.
The remaining requirement is $S_B\succeq0$.

For native sources, the estimates in Appendix~\ref{app:source} give
\[
 C\geq\frac{f_0}{4R_0HT},\qquad
 D\geq\frac{1-\delta^2}{4R_0T(1-H)}.
\]
Since $f_0<1-\delta^2$, their reciprocal sum is at most
$4R_0T[H+(1-H)]/f_0$. Thus
\begin{equation}
 \frac{CD}{C+D}\geq\frac{f_0}{4R_0T}
 >\frac1{5T}\geq\frac{1024}{5}>\frac{257}{2}.
 \label{eq:bobharmonic}
\end{equation}
Here $f_0/(4R_0)=140600049270784/564050540888577>1/5$.
The left side of \eqref{eq:bobharmonic} is below both $C$ and $D$.
Since $d=1$ on the native family, $S_B$ has positive diagonal entries
and determinant $CD-\tau_B(C+D)>0$.

Equations~\eqref{eq:bobparent}--\eqref{eq:boballocation} give the ten
fixed functions
\[
 \Phi_B=(E_\lambda,L_A,X,M_B,Y,U_B,J_B,cE_B,eE_B,
                          \sqrt{c^2-1/257}\,K_B)^T.
\]
Their matrix $Q_B(c)$ has diagonal
\[
 (\lambda,\beta,n+d/2,m,A,\alpha+d/2,
                   D-128d,C,C-257d/2,257d/2)
\]
and the single off-diagonal pair $(Q_B)_{6,8}=(Q_B)_{8,6}=257d/2$,
with indices starting at zero. It satisfies $S_c=\Phi_B^TQ_B(c)\Phi_B$.
Its $(8,6)$ principal block is $S_B$ plus $d/2$ on row~6, and all
other entries are nonnegative diagonal terms. At $d=0$ the same
matrix admits each event ray. The formulas therefore include $n=0$,
event additions and finite conic limits. They also apply when $b=0$
or $e=0$, covering the projective endpoints in this testing region.

Impose $Q_B(c)\succeq0$ together with the cap constraints of
Appendix~\ref{app:functions}. The resulting system has total order
$274+10+9=293$, with the nine source-sign blocks shared.
It admits every native source and proves positivity whenever
$|u|\leq9/256$ or $|v|\geq1/16$.
The unaddressed R1 testing region is $|u|>9/256$, $|v|<1/16$,
including Alice's projective endpoint in this strip.
The joined cone is closed: the cap cone is closed, and $Q_B(c)$ is
an affine matrix in the source coordinates with no additional variables.

The fifteen upper-triangular rational records of $Q_B$ are in
\path{bob-complement/bob_complement_order10.coo.json}; the adjoining
checker verifies every source coefficient in the polynomial identity.
\path{joined_certificate.json} specifies how this matrix and the cap
equations share their source coordinates.
